% =====================================================================
% セットアップ編 第2章 GitHubを使えるようにする
% 元原稿: chapters/02_github.tex（本編から分離）
% =====================================================================
\chapter{GitHubを使えるようにする}

この章では、\textbf{GitHubのアカウントを作り、自分のPCからGitHubにログインするまで}を扱います。ここまで済ませておくと、本編でGitHubに\cmd{push}（アップロード）するとき、認証で立ち止まらずに済みます。前章と同じく、\textbf{ここも一度やれば終わり}です。

なお、GitHubがどういうサービスなのかという説明は、本編『はじめてのバージョン管理』の第2章で行います。この章は\textbf{準備だけ}に集中します。

% ---------------------------------------------------------------------
\section{GitHubアカウントを作る}

まだアカウントがなければ作成します。

\begin{enumerate}[leftmargin=1.6em]
  \item \url{https://github.com/} を開き、\textbf{「Sign up」} をクリックします。
  \item メールアドレス・パスワード・ユーザー名を入力して進めます。
  \item 登録したメールアドレスに届く確認コードを入力すれば完了です。
\end{enumerate}

\begin{hintbox}
ここで登録する\textbf{メールアドレスは、前章の \cmd{git config -{}-global user.email} で設定したものに揃える}のがおすすめです。揃えておくと、コミットとGitHubアカウントが正しく結びつきます。
\end{hintbox}

% ---------------------------------------------------------------------
\section{GitHubにログインする（認証）}

GitHubに \cmd{push}（アップロード）するときは、「本当にあなたか」を確認する\textbf{認証}が必要です。本シリーズでは、いちばん手軽な \textbf{GitHub CLI（\cmd{gh}）でのログイン}を使います。これで一度ログインしておくと、\cmd{gh} がGitの認証を裏で肩代わりしてくれるため、\textbf{トークンを自分で発行しなくても \cmd{push} / \cmd{pull} がそのまま通ります}。

\begin{hintbox}
SSH鍵やトークン（PAT）を手動で使う方式もあります。頻繁に使うようになったら、『逆引きチートシート』の「認証の別方式（SSH / PAT）」を参照してください。まずは \cmd{gh} だけでOKです。
\end{hintbox}

\noindent\textbf{手順1：\cmd{gh} をインストールする}

\begin{warnbox}
\textbf{\cmd{gh} はGit本体には含まれていません。} 第1章でGitを入れても \cmd{gh} は入らないので、別途インストールが必要です。まず \cmd{gh -{}-version} を実行し、\cmd{command not found} なら下記で入れてください。
\end{warnbox}

\begin{codebox}[bash]
\begin{code}
# macOS（Homebrew）
brew install gh

# Windows（winget）
winget install --id GitHub.cli

# Linux（Debian / Ubuntu 系）
sudo apt install gh
\end{code}
\end{codebox}

インストール後、\cmd{gh --version} でバージョンが表示されれば成功です（うまくいかない場合は公式手順 \url{https://cli.github.com/} を参照）。

\noindent\textbf{手順2：ログインする}

\begin{codebox}[bash]
\begin{code}
gh auth login
\end{code}
\end{codebox}

対話形式で聞かれるので、次のように選びます。

\begin{center}
\begin{tabularx}{\linewidth}{@{}l X@{}}
\toprule
質問 & 選ぶもの \\
\midrule
Where do you use GitHub? & \textbf{GitHub.com} \\
preferred protocol for Git operations? & \textbf{HTTPS} \\
Authenticate Git with your GitHub credentials? & \textbf{Y} \\
How would you like to authenticate? & \textbf{Login with a web browser} \\
\bottomrule
\end{tabularx}
\end{center}

質問文や順序は \cmd{gh} のバージョンによって少し変わることがあります。上の表と完全に一致しなくても、\textbf{同じ意味の選択肢}（GitHub.com／HTTPS／ブラウザでログイン）を選べば問題ありません。

最後の選択でワンタイムコードが表示され、Enterを押すとブラウザが開きます。表示されたコードを入力し、\textbf{Continue} を押せば認証完了です。

\screenshot[0.62]{02_github/gh_auth_login.png}{GitHubログイン画面（Continueを押す）}
\screenshot[0.62]{02_github/gh_authorize_window.png}{GitHub認証画面（ワンタイムコードを入力）}

これで \cmd{gh} がGitの認証を肩代わりし、\textbf{以降 \cmd{git push} / \cmd{git pull} で認証を聞かれなくなります}。

\begin{codebox}[bash]
\begin{code}
# ログイン状態を確認する
gh auth status
\end{code}
\end{codebox}

\begin{warnbox}
トークンの有効期限が切れたり、GitHub側でアクセスを取り消したりすると、認証を聞かれるようになります。そのときは \cmd{gh auth login} で入り直してください。
\end{warnbox}


% ---------------------------------------------------------------------
\section{セットアップ完了チェック}

ここまでで準備は終わりです。次の4つを実行して、\textbf{すべてエラーなく表示されれば完了}です。本編『はじめてのバージョン管理』へ進んでください。

\begin{codebox}[bash]
\begin{code}
git --version              # Gitのバージョンが表示される
git config --global --list # user.name と user.email が表示される
gh --version               # ghのバージョンが表示される
gh auth status             # Logged in to github.com と表示される
\end{code}
\end{codebox}

\begin{center}
\begin{tabularx}{\linewidth}{@{}l X@{}}
\toprule
確認すること & うまくいかないときは \\
\midrule
Gitが入っているか & 第1章「Gitのインストール」へ戻る \\
名前・メールが登録されているか & 第1章「最初の設定」へ戻る \\
\cmd{gh} が入っているか & 本章「GitHubにログインする」の手順1へ戻る \\
GitHubにログインできているか & \cmd{gh auth login} をもう一度実行する \\
\bottomrule
\end{tabularx}
\end{center}

インストールしたはずなのに \cmd{command not found} と出るときは、\textbf{ターミナルをいったん閉じて開き直して}から、もう一度試してください。インストール直後は、設定がターミナルに反映されていないことがあります。
